\section{Introdución}
% Relatividade en 4D (acción EH?). Necesidade de D>4. Limitación EH en D>4. Necesidade do teorema de Lovelock. Contexto moderno: gravidade en altas dimensións, \textbf{supergravidade} e teorías gauge.

%Un dos principais problemas históricos da ciencia é a creación dunha \textit{teoría do todo}, hoxe en día a maior traba para a realización deste soño é a cuantificación da interacción gravitacional. A vía que se atopou como natural consiste en converter a relatividade xeral de Einstein nunha teoría gauge de Yang-Mills vía renormalización, tras máis de 70 anos de traballo científico atopouse que isto é imposible para unha teoría gravitacional en $3+1$ dimensións, aínda que si para o caso de $2+1$ dimensións \cite{2+1grav0,2+1grav1}. Semella que o camiño a seguir é crear novas teorías de gravidade de dimensións maiores a $3+1$ (teorías de supergravidade?) o camiño abriuno David Lovelock no seu traballo de 1971 \cite{Lovelock-1971} que baixo o teorema que leva o seu nome permite elaborar certos tipos de accións (accións de Chern-Simons?) para gravidade a altas dimensións.

%O noso punto de partida será a formulación orixinal de Einstein de 1915 coa acción de Einstein-Hilbert que xeneralizaremos á teoría de Cartan, probaremos que acaba conducindo a unha formulación semellante. A relatividade xeral queda reducida a unha teoría efectiva. principio de equivalencia de Einstein

Un dos principais problemas históricos da ciencia é a creación dunha \textit{teoría do todo}. Neste contexto, a gravidade presenta dificultades especiais, xa que a súa descrición clásica mediante a Relatividade Xeral de Einstein non se integra de forma natural no marco das teorías cuánticas de campos. En particular, os intentos de tratar a gravidade como unha teoría gauge análoga ás teorías de Yang–Mills conducen a problemas de non renormalizabilidade en $3+1$ dimensións, incluso se proba que é imposible, aínda que para o caso de $2+1$ dimensións \cite{2+1grav0,2+1grav1} si que é posible.

Semella que o camiño a seguir é crear novas teorías de gravidade de dimensións superiores. Así, David Lovelock, no seu traballo de 1971 \cite{Lovelock-1971} proporciona un resultado fundamental baixo o teorema que leva o seu nome que permite elaborar formas máis xerais da acción gravitacional
\begin{align}
    \mathcal{S}_{EH}=\int \dd^D x \sqrt{\abs g} \frac{1}{2\kappa}R \hspace{1.5em} \text{onde} \hspace{1em} \kappa=8\pi G.
\end{align}
O noso punto de partida será a formulación orixinal de Einstein de 1915 coa acción de Einstein-Hilbert (1) que xeneralizaremos á teoría de Cartan, ademais de amosar outras vías, até introducir naturalmente as teorías de Lovelock, entendidas como extensións consistentes da gravidade en dimensións arbitrarias.

\section{Teoría de variedades}
Antes de poder describir as teorías de Lovelock debemos construír unha adecuada estrutura matemática na que fundamentar a nosa teoría. O primeiro paso será destripar a nosa mesa de traballo que chamamos espazotempo, diremos que o espazotempo é unha $\mathcal{M}$ \textit{variedade pseudoriemanniana afinmente conectada}. Daremos por coñecidos as nocións fundamentais e a álxebra de tensores, tales como cambios de índice, vectores co- e contravariantes como tamén os cambios xenerais de coordenadas (CXC), trataremos todo nesta sección do xeito máis breve e cualitativo posible.

\subsection{Variedades diferenciais}
O nivel de estrutura máis baixo dende o que podemos estender a noción de variedade até que inclúa todos os anteriores palabros é a variedade topolóxica que é un tipo de estrutura alxébirca ao que se lle asocia un homeomorfismo \cite{masa-2020}
\begin{adjustwidth}{1cm}{0.8cm}\vspace{0.4em}
    \textit{Sexan $\mathcal{M},\mathcal{N}\subset\mathbb{R}^p,\mathbb{R}^q$ dous espazos. Diremos que unha aplicación continua e con inversa continua $h:\mathcal{M}\longrightarrow\mathcal{N}$ é un \textbf{homemorfismos}. Se existe $h$ diremos que $\mathcal{M}$ e $\mathcal{N}$ son homeomorfos, $\mathcal{M}\approx\mathcal{N}$.}\vspace{0.4em}
\end{adjustwidth}
No entanto, na física, é moito máis interesante estudar un tipo concreto de variedades topolóxicas sobre a que si podemos realizar ``cálculos'': as variedades diferenciais; estas caracterízanse por ter asociado un difeomorfismo, que non é máis ca unha restrición da existencia de derivadas sobre os homeomofismo e a inversa. Dicimos que aquelas propiedades que se conservan baixo un homeomorfismo son propiedades topolóxicas, as que se conserven baixo os difeomorfismo diremos que son xeométricas. Entre estas propiedades destacamos a metricidade e a conexidade.

\subsubsection{Metricidade}
A métrica $g\indices{_\mu_\nu}$ é o obxecto que nos permite medir distancias dentro da nosa variedade, trátase dun campo tensorial simétrico e non dexenerado. Dicimos que unha \textit{variedade diferencial é de Riemann} cando esta presenta unha estrutura de métrica definida positiva, denotámolo por ($\mathcal{M},\textbf{g}$). Nós utilizamos un tipo algo máis xeral de variedade, aquelas denominadas \textit{pseudoriemannianas}, estas son aquelas cuxa métrica non é necesariamente definida positiva. Logo, para unha variedade desta clase definimos a signatura da métrica como a tupla ($p,q$) onde $p$ representa a cantidade de autovalores de valor $-1$ e $q$ a de $+1$. Notemos que a propiedade de simetría permite sempre diagonalizar a métrica, o que quere dicir que podemos construír unha forma canónica, para o primeiro caso teríamos unha métrica canónica euclídea $\delta\indices{_\mu_\nu}$ e para o segundo unha métrica canónica de Minkowski $\eta\indices{_\mu_\nu}$ de signatura ($p,q$).

A métrica permítenos definir tamén o produto interno entre vectores que á vez permítenos medir a norma destes. No caso dunha variedade pseudoriemanniana dicimos que esta norma non ten signo definido e segundo se é negativa, nula ou positiva, dicimos que é o vector é \textit{time-like}, \textit{light-like} ou \textit{space-like}, respectivamente.

\subsubsection{Afinidade}
Se ben poderíase pensar que a métrica é suficiente para estudar a xeometría da variedade por extensión euclidiana, o certo é que atopamos que aquelas variedades máis xenerais ``conectan'' os seus puntos de xeitos non moi intuitivos. Unha variedade afinmente conectada é aquela que presenta unha conexión $\Gamma$ (o que nos di ``como se conecta'' o espazo) e denotámola como ($\mathcal{M},\Gamma$). Isto vese reflectido cando pretendemos transportar un vector dun punto a outro da variedade\footnote{Isto vese explicitamente na derivada covariante, revísese \cite{notasclase,doCarmo,lovelock-1989}. Constatar só que para poder definir correctamente unha derivada covariante a conexión é imprescindible.}, onde as propiedades só se conservan na translación se a estrutura é afín; i.e. transporte autoparalelo, que nos dá unha noción de ``recto'' no sentido de un vector mantéñase paralelo a si mesmo ao longo do transporte.

\subsection{Xeodésicas sobre unha variedade}
A priori, a conexión e a métrica son obxectos que nada di un do outro e en consecuencia poderíamos pensar en dous tipos de traxectorias óptimas distintas sobre a variedade: unha que sexa a máis corta en canto a métrica e outra que sexa a máis recta en canto o transporte é autoparalelo. A primeira clase é a traxectoria tipicamente física que supón unha curva crítica para unha acción, como é o caso da $\mathcal{S}_{EH}$, estas curvas dependen dun obxecto denominado \textit{símbolos de Christoffel} $\{ \indices{^\mu_\alpha_\beta}\}$ e rexidas pola ecuación da xeodésica métrica (2). O segundo tipo, que denominamos xeodésica afín (3), depende da conexión e define a traxectoria máis recta ou autoparalela para un vector. 
\begin{align}
    \ddot{x}\indices{^\mu} + \{ \indices{^\mu_\alpha_\beta}\}\dot{x}\indices{^\alpha}\dot{x}\indices{^\beta} = 0\\
    \ddot{x}\indices{^\mu} + \Gamma\indices{^\mu_\alpha_\beta}\dot{x}\indices{^\alpha}\dot{x}\indices{^\beta} = 0
\end{align}
Os símbolos de Christoffel veñen completamente determinados pola métrica como se indica en (4). Isto é así porque a métrica é a única que media en xeral na definición dunha densidade lagranxiana, utilizando Euler-Lagrange chegaríamos a
\begin{align}
    \{ \indices{^\mu_\alpha_\beta}\} \equiv \levicivita
\end{align}

\subsection{Curvatura e torsión}
A partir da conexión, e unicamente desta, podemos construír dous tensores que miden dun modo moito máis directo ou intuitivo o comportamento xeométrico da variedade e que nos permitirá distinguir entre distintos espazotempos. Nun primeiro lugar construímos o tensor de torsión $T\indices{^\sigma_\mu_\nu}$ que cuantifica como rotan os vectores ao longo do transporte, digamos que mide o ``antiparalelismo''\footnote{A interpretación máis usual é algo como ``\textit{mide en que medida falla o peche do paralelogramo que formaría o transporte de dous vectores entre eles''}.}. É a parte antisimétrica da conexión
\begin{align}
    T\indices{^\sigma_\mu_\nu} = \Gamma\indices{^\sigma_\mu_\nu} - \Gamma\indices{^\sigma_\nu_\mu},
\end{align}
notemos que a torsión é nula se a conexión é simétrica nos índices inferiores, se a variedade non ten torsión dise libre de torsión.
Pola contra, o tensor de Riemann $R\indices{^\mu_\nu_\alpha_\beta}$ contén toda a información sobre a curvatura xeométrica da variedade
\begin{align}
    R\indices{^\mu_\nu_\alpha_\beta} = \riemann ,
\end{align}
se todas as compoñentes son nulas dise que a variedade é plana e deberímos ter un espazo difeomorfo ao de Minkowski. Ademais este tensor pode ter tres contraccións independentes posibles moi importantes
\begin{align}
    R\indices{^{(1)}_\mu_\nu}&:= R\indices{^\sigma_\mu_\sigma_\nu} = \partial_\sigma \Gamma\indices{^\sigma_\mu_\nu} - \partial_\nu \Gamma\indices{^\sigma_\mu_\sigma} + \Gamma\indices{^\sigma_\lambda_\sigma}\Gamma\indices{^\lambda_\mu_\nu} - \Gamma\indices{^\sigma_\lambda_\nu}\Gamma\indices{^\lambda_\mu_\sigma}\\
    R\indices{^{(2)}_\mu_\nu}&:= R\indices{_\mu^\sigma_\sigma_\nu} = g\indices{_\rho_\mu}g\indices{^\sigma^\varsigma}\left( \partial_\sigma \Gamma\indices{^\rho_\varsigma_\nu} - \partial_\nu \Gamma\indices{^\rho_\varsigma_\sigma} + \Gamma\indices{^\rho_\lambda_\sigma}\Gamma\indices{^\lambda_\varsigma_\nu} - \Gamma\indices{^\rho_\lambda_\nu}\Gamma\indices{^\lambda_\varsigma_\sigma} \right)\\
    R\indices{^{(3)}_\mu_\nu}&:= R\indices{^\sigma_\sigma_\mu_\nu} = \partial_\mu \Gamma\indices{^\sigma_\nu_\sigma} - \partial_\nu \Gamma\indices{^\sigma_\mu_\sigma} 
\end{align}
que coa métrica (notemos que para $R\indices{^{(2)}_\mu_\nu}$ xa foi preciso empregala) podemos calcular a traza destes tensores
\begin{align}
    R\equiv R^{(1)} = -R^{(2)}\hspace{2em} \text{e} \hspace{2em} R^{(3)} = 0.
\end{align}
Só un destes escalares é non trivial, é o denominado \textit{escalar de curvatura} que constrúe o tensor de Einstein e aparece na acción de (1), por outro lado o tensor $R\indices{^{(1)}_\mu_\nu}$ é un tanto especial e denominámolo tensor de Ricci.

\subsection{Teoría de Riemann-Cartan}
Coas primeiras definicións xa introducimos o que é unha variedade pseudoriemanniana afinmente conectada ($\mathcal{M},g,\Gamma$), como esixe o noso espazotempo. No entanto, imos restrinxir un pouco máis a nosa variedade para conseguir unha contracción de Riemann única e non haxa ambigüidade na definición de curvatura. Até agora nada dixemos sobre a relación entre métrica e conexión, no entanto, podemos avaliar a compatibilidade da métrica mediante o tensor de metricidade\footnote{Este tensor defínese mediante a derivada covariante da métrica, fagamos unha pequena aclaración sobre notación. En xeral referirémonos a esta cun ``$;\mu$'' cando non lle esixamos ningunha clase de restrición (conexión arbitraria), no entanto, para unha conexión concreta usaremos $\nabla_\mu$ cun acento que represente á conexión precisa, por exemplo $\breve\nabla_\mu$ para Levi-Civita.}
\begin{align}
    Q\indices{_\lambda_\mu_\nu} = -g\indices{_\mu_\nu_{;\lambda}},
\end{align}
se este anúlase dicimos que a conexión é compatible coa métrica; se escollemos unha conexión baixo a que a variedade satisfaga esta condición dicimos que estamos ante unha teoría de Riemann-Cartan. Neste caso recuperamos as propiedades alxébricas (simetrías) típicas do tensor de Riemann e é inmediato que a única contracción posible é o tensor de Ricci. A verdadeira razón de esixirlle ao noso espazotempo que sexa unha teoría de Riemann-Cartan é que a compatibilidade coa métrica converte ao transporte paralelo ao longo dunha traxectoria nunha isometría, i.e., o produto escalar, logo, tamén a norma dun vector, é invariante ao longo da curva; este resultado é coñecido como \textit{Lema de Ricci} \cite{lovelock-1989}. Esta é unha propiedade necesaria, pois non sería fisicamente admisible que un vector tipo luz convértase nun vector doutra clase segundo o punto onde se atope: non teríamos correspondencia fenomenolóxica. 

\paragraph{Teorema fundamental da xeometría riemanniana}
\begin{adjustwidth}{1cm}{0.8cm}
    \textit{Sexa unha variedade ($\mathcal{M},g,\Gamma$) existe unha única conexión que é compatible coa métrica e simétrica que denominamos \textit{conexión de Levi-Civita}, $\breve\Gamma$.} \cite{doCarmo}\vspace{0.4em}
\end{adjustwidth}


Esta conexión é moi especial, non é só que sexa un bo candidato como conexión para un espazotempo, senón que ademais é a única posible cuxa variedade é libre de torsión (por ser simétrica), ademais a súa expresión analítica é 
\begin{align}
      \breve\Gamma\indices{^\mu_\alpha_\beta}= \levicivita \equiv \{ \indices{^\mu_\alpha_\beta}\},
\end{align}
os símbolos de Christoffel coinciden coa conexión de Levi-Civita! O que se traduce que as traxectorias críticas (2) e autoparalelas (3) son as mesmas. Veremos na sección 4.2 que esta non é a única conexión que verifica que as xeodésicas coincidan, no entanto si que simplifican moito os cálculos. Tamén se pode demostrar que na teoría de Riemann-Cartan unha conexión calquera pode ser expresada como
\begin{align}
    \Gamma\indices{^\mu_\alpha_\beta} &= \breve\Gamma\indices{^\mu_\alpha_\beta} + K\indices{^\mu_\alpha_\beta} + L\indices{^\mu_\alpha_\beta}\\
    \nonumber K\indices{^\mu_\alpha_\beta} &= \frac{1}{2}g\indices{^\mu^\sigma}\left(T\indices{_\sigma_\alpha_\beta} + T \indices{_\nu_\sigma_\beta} - T\indices{_\alpha_\beta_\sigma}\right)\\
    \nonumber L\indices{^\mu_\alpha_\beta} &= \frac{1}{2}g\indices{^\mu^\sigma}\left(Q\indices{_\sigma_\alpha_\beta} + Q \indices{_\nu_\sigma_\beta} - Q\indices{_\alpha_\beta_\sigma}\right)
\end{align}
onde $K\indices{^\mu_\alpha_\beta}$ é o tensor de contorsión. A conexión depende só da torsión e da metricidade, evidenciando o teorema anterior.

\section{Formas diferenciais e xeometría exterior}
Entre 1929 e 1932 Einstein e Cartan sostiveron un debate por correspondencia postal \cite{CH_supergravities} pola veracidade da conexión de Levi-Civita como a lexítima da relatividade xeral. Efectivamente esta satisfai as condicións para que a nosa teoría sexa física e ademais é especialmente doada para reproducir cálculos (este é o que poderíamos chamar punto de vista de Einstein). No entanto, Cartan argumentaba que a maior xeneralidade obteríase separando a métrica da conexión; hai dous xeitos de lograr isto, cunha conexión xeral que se denomina conexión de Palatini \cite{Lovelock_metricformulation,Bernal_2017}, ou, no lugar de traballar sobre a variedade, mediante o estudo do fibrado tanxente, o que se pode chamar punto de vista de Cartan que agora expoñemos.
\subsection{Álxebra exterior de formas diferencias}
Antes de falar \textit{per sé} da formulación de Cartan estudaremos uns obxectos previos que se denominan formas diferenciais que son imprescindibles para traballar no fibrado tanxente, de feito unha 1-forma será precisamente o vector dual (ou vector contravariante) dun vector do espazo tanxente $\tpm$, i.e., vive no espazo cotanxente $\cotpm$.
%\subsubsection{\textit{p}-formas e produto exterior}
% \subsubsection{Producto \textit{wedge}}
% \subsubsection{Derivada exterior}

Tomemos un vector $A_j$ dun campo covariante sobre unha variedade diferencial, cun sistema local de coordenadas (base holónoma), para un desprazamento arbitrario de compoñentes $\dd x^j$ podemos escribir a suma
\begin{align}
    \omega = A_j\dd x^j
\end{align}
que é unha 1-forma\footnote{Estas aparecen en infinidade de situacións da física baixo o nome de pfaffiana aínda que non se tratan xeometricamente, por exemplo nas ecuacións de Gibbs da termodinámica.}, notemos que esta é un escalar que satisfai trivialmente condicións de linearidade. Se no lugar dun vector escollemos un tensor, por exemplo, (1,1), obteríamos unha 1-forma contravariante.\\
Se agora quixeramos operar entre dúas 1-formas necesitaremos introducir un produto externo da estrutura alxébrica da variedade ao que denominamos produto \textit{wedge} e se denota como $\omega\wedge\pi$ onde $\omega$ e $\pi$ son 1-formas. O seu produto denominarémolo 2-formas (obviamente estas non pertence á variedade) e satisfai as propiedades de asociación e distribución usuais da multiplicación, no entanto este é anticomutativo
\begin{align}
    \omega\wedge\pi = -\pi\wedge\omega,\hspace{1.2em}\text{en particular} \hspace{1em} \dd x^j\wedge\dd x^h = - \dd x^h\wedge\dd x^j .
\end{align}
Con este produto definido podemos chegar case inmediatamente por indución á definición de p-formas como\footnote{Nesta ecuación poñemos un expoñente na forma para sinalar a súa orde, no sucesivo referirémonos a $\omega$ como unha p-forma e a $\pi$ como unha q-forma sen poñer o superíndice.}
\begin{align}
    \omega^p = A\indices{_j_h}\dd x^j\dd x^h\longrightarrow \omega^p = A\indices{_{j_1}_{...}_{j_p}}\dd x^{j_1}\wedge ... \wedge \dd x^{j_p}
\end{align}
e por tanto chegamos a probar que a anterior regra anticonmutativa é 
\begin{align}
    \omega\wedge\pi = (-1)^{pq}\pi\wedge\omega .
\end{align}
O noso estudo limítase ás 1 e 2-formas, no entanto, daremos estos primeiros resultados en xeral para p-formas que son trivialmente particularizados. 
% podería introducir aquí a construción linear e os conceptos de formas lineares independentes e que unha combinación linear \wedge\omega = 0
Agora introduciremos a derivada exterior das p-formas que necesitaremos para definir un novo tipo de derivada covariante, que así enunciada de corrido coas súas propiedades esta é
\begin{align}
    \Omega\equiv\dd\omega &= \diffp{A\indices{_{j_1}_{...}_{j_p}}}{x^k}\dd x^k\wedge \dd x^{j_1}\overset{\wedge}{...}\ \dd x^{j_p} = (-1)^p\diffp{A\indices{_{j_1}_{...}_{j_p}}}{x^l}\dd x^{j_1}\overset{\wedge}{...}\  \dd x^{j_p}\wedge\dd x^l,\\
    \dd(\omega +\pi) &= \dd\omega + \dd\pi\hspace{7.21em}\text{(Distributiva respecto da suma)},\\
    \dd(\omega\wedge\pi) &= \dd\omega\wedge\pi + (-1)^p\omega\wedge \dd\pi\hspace{1.2em}\text{(Distributiva respecto do wedge),}
\end{align}
notemos que a derivada dunha p-forma, $\Omega$, é unha (p+1)-forma.\\
Podemos chegar inmediatamente a un resultado moi importante no análise de formas diferencias mediante a antisimetría da definición (18), que xorde da anticonmutatividade (17), se pretendemos calcular a derivada segunda, atoparemos que esta debe anularse,
\begin{align}
    \dd^2\omega = \dd(\dd\omega) = 0,
\end{align}
xa que a derivada parcial é conmutativa mentras que as formas elementais non (18), este importantísimo resultado ao que recorreremos multitude veces é coñecido como \textit{lema de Poincaré}. Anecdoticamente, a partir deste pódese chegar ao teorema de Stokes máis xeral do cálculo, $\int_{\partial\mathcal{N}}\omega = \int_\mathcal{N}\dd\omega$.

% Poincaré's Lemma
% a derivada exterior leva inmediatamente ao teorema de Stokes (comentalo sen máis) 

% 3.2. Integración y dualidad
% Operador Hodge
% Integración sobre variedades
% Teorema de Stokes

\subsection{Formulación dende o fibrado tanxente}
%\subsubsection{Espazo tanxente sobre a variedade}
Sexa unha variedade diferencial D-dimensional $\mathcal{M}$, como o noso espazotempo, podemos construír unha aproximación plana inmediata para todo punto $P$ da variedade que denominamos espazo tanxente ($\tpm$), i.e., localmente toda variedade é homeomorfa ao espazo de Minkowski. Ao conxunto de todos os espazo tanxentes denominámolo fibrado tanxente $T\mathcal{M}$. Como para calquer espacio vectorial podemos construír o seu dual que denominamos espazo cotanxente $\cotpm$, así distinguimos entre vectores covariantes e contravariantes como usualmente. O convenio de índices será o seguinte: para representar a base holónoma (a propia da variedade) continuaremos utilizando índices gregos, para a base ortonormal na tanxente usamos índices latinos. 

\subsubsection{Vielbein}
Para poder escribir toda a relatividade en linguaxe do fibrado tanxente necesitamos poder movernos entre coordenadas locais (holónomas) e as ortonormais do tanxente, este será un isomorfismo da forma $\textbf{e}:\mathcal{M}\longrightarrow T\mathcal{M}$ que coincide coa matriz jacobiana usual. De feito, acción sobre coordenadas separadas infinitesimalmente cífrase por definición como
\begin{align}
    \diffp{z^a}{x^\mu}=\viel(x)\Rightarrow dz^a = \viel(x)\dd x^\mu.
\end{align}
A familia \{$\viel\ |\ a=1,...,D$\} é o que chamamos \textit{vielbein}. A partir da definición anterior podemos decatarnos inmediatamente dalgunhas propiedades fundamentais deste obxecto como
\begin{align}
    \viel e\indices{^\mu_b}= \delta^a_b\ ,\hspace{1.2em} e\indices{^\mu_a} e\indices{^a_\nu} =\delta^\mu_\nu\ ;\hspace{2em} &\text{(ortonormalidade do índice)}\\
    T\indices{^\mu^\nu}= \vieli T\indices{^a^\nu} = \vieli\bieli T\indices{^a^b}\ ;\hspace{2em} &\text{(transformación entre espazos de tensores)}
\end{align}
a transformación de tensores é facilmente xeneralizable para tensores (p,q). Notemos que podemos reescribir a métrica co vielbein, a partir do elemento de arco para o espazo tanxente atopamos que 
\begin{align}
    \dd s^2 =\eta\indices{_a_b}\dd z^a\dd z^b = \eta\indices{_a_b}\viel(x)\biel(x)\dd x^\mu\dd x^\nu\Rightarrow g\indices{_\mu_\nu}(x) = \viel(x)\biel(x)\eta\indices{_a_b}. 
\end{align}
Posto que o vielbein determina a métrica é intuitivo pensar que todas as propiedades métricas do espazotempo están contidas no vielbein. O contrario non é certo pois pódense construír moitos vielbein que reproduzan a mesma métrica, isto débese a que hai moitas eleccións de bases ortonormais no tanxente, de feito notemos que podemos transformar o vielbein baixo transformacións locais de Lorentz (TLL) ${\mathbf{\Lambda}(x)\in SO(D-1,1)}$ tal que estas pasan desapercibidas para a variedade $\mathcal{M}$ (porque a métrica $g\indices{_\mu_\nu}$ é invariante baixo TLL)
\begin{align}
    \viel(x)\longrightarrow e\indices{^{a'}_\mu}(x) = \Lambda\indices{^{a'}_a}(x)\viel(x).
\end{align}
Unha consecuencia inmediata disto é que, en xeral, os índices latinos transforman baixo TLL mentres que os índices gregos transforman como a métrica, baixo CXC, i.e., Diff($\mathcal{M}$).

\subsubsection{Conexión de Lorentz}
    % só será necesario para campos de spin racional, como o de Dirac
O primeiro problema que debemos solucionar na formulación de Cartan para obter unha teoría equivalente á usual da relatividade é definir un operador diferencial que transforme adecuadamente baixo TLL, que son as que afectan aos obxectos (tensores) que vivan no tanxente: unha derivada covariante. Notemos a simetría do problema para os CXC sobre a variedade. Necesitamos logo un tipo de conexión que se chama conexión de Lorentz ou conexión de spin $\omega\indices{^a_b_\mu}$, así a nosa derivada covariante para un campo $\phi^a(x)$ 
\begin{align}
    \mathcal{D}_\mu \phi^a = \partial_\mu\phi^a + \omega\indices{^a_b_\mu}\phi^b.
\end{align}
Esta definición require que a conexión transforme como
\begin{align}
    \omega\indices{^a_b_\mu}\longrightarrow \omega\indices{^{a'}_{b'}_\mu} = \Lambda\indices{^{a'}_a}\Lambda\indices{_{b'}^b} \omega\indices{^a_b_\mu} + \Lambda\indices{^{a'}_a}\partial_\mu\Lambda\indices{_{b'}^b}.
\end{align}
Ao igual que pasaba con $\viel$ e a métrica, as propiedades afíns da variedade están codificadas en $\omega\indices{^a_b_\mu}$ que é independente da métrica, ao igual que cunha conexión xeral de Palatini. Esta conexión de spin é análoga á conexión afín (asimesmo pode definir un transporte paralelo no fibrado), malia ser obxectos independentes por vivir en espazos diferentes, estos poderían ser relacionados polo isomorfismo entre $\mathcal{M}$ e $T\mathcal{M}$, $\viel$.

\paragraph{Vielbein e conexión de spin como 1-formas.}
% se sobra espazo aquí introduciría un comentario sobre a equivaleencia de vielbein coa conexion afin e a torsion e curvatura usuais para facer un apéndice
Todas as propiedades xeométricas de $\mathcal{M}$ poden ser expresadas mediante dúas 1-formas xunto cos seus productos e derivadas exteriores. 
Estas 1-formas necesarias constrúense mediante a supresión da dependencia de coordenadas holónomas, i.e., contraendo o índice grego, das nosas estudadas estruturas do fibrado: o vielbein e a conexión de Lorentz. Logo, as seguintes construccións
\begin{align}
    e^a&=\viel\dd x^\mu\\
    \omega\indices{^a_b} &=\omega\indices{^a_b_\mu}\dd x^\mu 
\end{align}
son invariantes baixo CXC, o que se di libre de coordenadas. Por construción todas os produtos e derivadas exteriores tamén o son, de feito a derivada covariante anterior transfórmase no operador derivada covariante exterior para un campo tensorial
\begin{align}
    \mathcal{D} = \dd +\omega\indices{^a_b}\wedge
\end{align}
onde `$\dd$' representa a derivada exterior.
\subsubsection{Curvatura e torsión como 2-formas}
%   Ecuacións de estrutura de Cartan p165\cite{lovelock-1989}
Outra consecuencia do lema de Poincaré (21) é que o cadrado do operardor (31) non é un operador diferencial, no entanto, si que é un operador alxébrico que é coñecido como a curvatura 2-forma. Desenvolvendo (31) cun campo tensorial auxiliar $\phi^a$
\begin{align}
    \hspace{-2em}\nonumber\mathcal{D}^2\phi^a = \mathcal{D}(\mathcal{D}\phi^a) &= \dd(\mathcal{D}\phi^a)+\omega\indices{^a_b}\wedge\mathcal{D}\phi^b = \dd(\dd\phi^a +\omega\indices{^a_b}\wedge\phi^b) + \omega\indices{^a_b}\wedge(\dd\phi^b +\omega\indices{^b_c}\wedge\phi^c)\\
    \nonumber &= \dd^2\phi^a + \dd\omega\indices{^a_b}\wedge\phi^b - \omega\indices{^a_b}\wedge\dd\phi^b + \omega\indices{^a_b}\wedge\dd\phi^b + \omega\indices{^a_b}\wedge\omega\indices{^b_c}\wedge\phi^c\\
     &= \dd\omega\indices{^a_c}\wedge\phi^c + \omega\indices{^a_b}\wedge\omega\indices{^b_c}\wedge\phi^c = (\dd\omega\indices{^a_c} + \omega\indices{^a_b}\wedge\omega\indices{^b_c})\wedge\phi^c \equiv \mathcal{R}\indices{^a_b}\wedge\phi^b.
\end{align}
A expresión (32) define un tensor (1,1) no fibrado tanxente que é unha 2-forma, esta tamén recibe o nome de \textit{primeira ecuación de estrutura de Cartan}. Esta curvatura pódese relacionar directamente co tensor de Riemann usual (6) cunha serie de asuncións que non explicaremos aquí, %(véxase o apéndice A)
(véxase \cite{lovelock-1989,EinsteinsVielben,clave} para unha explicación pormenorizada)
tal que
\begin{align}
    \mathcal{R}\indices{^a_b} = \dd\omega\indices{^a_b} + \omega\indices{^a_l}\wedge\omega\indices{^l_b} = \frac{1}{2}\viel\bieli R\indices{^\mu_\nu_\rho_\sigma}\dd x^\rho\wedge\dd x^\sigma = \eta\indices{^a^c}e\indices{^\mu_c}\bieli R\indices{_\mu_\nu}.
\end{align}
A outra 2-forma que construímos a partir das nosas 1-formas é empregando a derivada exterior covariante da 1-forma (29) 
\begin{align}
    \mathcal{T}^a = \dd e^a + \omega\indices{^a_b}\wedge e^b
\end{align}
que é un tensor (1,0) no fibrado tanxente que é a 2-forma da torsión, esta tamén recibe o nome de \textit{segunda ecuación de estrutura de Cartan}. Ao igual que antes, esta 2-forma de torsión podemos relacionala co tensor de torsión usual (5) tal que
\begin{align}
    \mathcal{T}\indices{^a} =\dd e^a + \omega\indices{^a_b}\wedge e^b = \frac{1}{2}\viel T\indices{^\mu_\rho_\sigma}\dd x^\rho\wedge\dd x^\sigma.
\end{align}
A partir destas 2-formas poderíamos pensar en seguir construíndo formas de maior orde mediante a derivada covariante exterior, no entanto de aplicala sobre (32) e (34) obteríamos
\begin{align}
    \mathcal{D}\mathcal{R}\indices{^a_b} =\dd \mathcal{R}\indices{^a_b}  + \omega\indices{^a_c}\wedge\mathcal{R}\indices{^c_b} - \omega\indices{^c_b}\wedge\mathcal{R}\indices{^a_c} \\
    \mathcal{D}\mathcal{T}\indices{^a} = \dd\mathcal{T}\indices{^a} + \omega\indices{^a_b}\wedge T^b = \mathcal{R}\indices{^a_b}\wedge e^b
\end{align}
é dicir, non podemos construír máis estructuras non triviais no formalismo de Cartan. As expresións (36) e (37) reciben o nome de \textit{identidades de Bianchi} e \textit{identidade cíclica} para a tanxente, respectivamente. A proba é trivial, véxase, p.ex.,
\cite{lovelock-1989}.
% 3.4. Ecuaciones de Cartan
% Primera ecuación estructural
% Segunda ecuación estructural


\section{Dinámica lagranxiana na relatividade xeral}
Unha das ferramentas máis poderosas e elegantes da física son, sen lugar a dúbida, os principios variacionais. A partir dun funcional de acción adecuado pódese establecer un principio de mínima acción do mesmo que nos permita recuperar toda a dinámica asociada a esta, no caso dunha acción derivada por unha lagranxiana dispomos das ecuacións de Euler-Lagrange coas que, en principio, deberíamos recuperar as ecuacións dos campos asociados á lagranxiana, en xeral
\begin{align}
    \delta\mathcal{S} = 0=\delta\int\mathcal{L}(q,\dot q)\sqrt{\abs g}\dd^D x\Rightarrow \diffp {\sqrt{\abs g}\mathcal L}{q}-\diff*{\diffp{\sqrt{\abs g}\mathcal L}{\dot{q}}[]}\tau = 0
\end{align}
Para o caso da relatividade xeral usual, a acción que podemos enunciar é a acción de Einstein-Hilbert (1) coa que recuperamos as ecuacións de Einstein no baleiro, no entanto, para recuperar o caso con materia basta con asociarlle unha lagranxiana material do tipo $\mathcal{L}_{\text{grav}}(\chi,g,\Gamma)$\footnote{Sobreentenderase que as lagranxianas dependen tamén das primeiras derivadas.} onde $\chi$ refire ao campo físico preciso, probaremos esta equivalencia no seguinte apartado. Se asumimos que a gravidade vén completamente determinada polas propiedades xeométricas da variedade a lagranxiana máis xeral que podemos enunciar é $\mathcal{L}_{\text{grav}}(g,\Gamma)$, obviamente este é o caso para (1), pero esta non é a máis xeral, as lagranxianas máis xerais serán as de Lovelock.

\subsection{Formalismo métrico}
O formalismo métrico, ou de segunda orde, é o que se corresponde co punto de vista de Einstein polo que asume que a conexión lexítima é a de Levi-Civita (12). Isto quere dicir, como xa vimos, que a nosa gravidade estará completamente determinada por tan só a métrica e as súas derivadas primeiras e segundas\footnote{Coidemos que as lagranxianas, en xeral, escríbense só coas primeiras derivadas, no entanto, como a conexión $\breve\Gamma\sim\partial g$ a súa derivada $\partial\breve\Gamma\sim\partial^2 g$.} e ademais será libre de torsión. 
Probemos agora polo principio de mínima acción que podemos reconstruír as ecuacións de Einstein variando a métrica\footnote{Asumiremos que $\delta\breve R\indices{_\mu_\nu} = 0$, véxase a demostración en \cite{clave}}.
  
\begin{align}
    \hspace{-3em}\nonumber\mathcal{S}=\mathcal{S}_{EH} + \mathcal{S}_{\text{material}} = \int\dd^D x \sqrt{\abs g} \left[ \frac{1}{2\kappa}\breve R +\mathcal{L}_{\text{material}}(\chi, g)\right] = \int\dd^D x \sqrt{\abs g} \left[ \frac{1}{2\kappa}g\indices{^\mu^\nu}
    \breve R\indices{_\mu_\nu} +\mathcal{L}_{\text{material}}(\chi, g)\right]
\end{align}
\begin{align*}
    \delta\mathcal{S} = 0 &= \int\dd^D x  \left[ \frac{1}{2\kappa}\delta\left(\sqrt{\abs g}g\indices{^\mu^\nu}
    \breve R\indices{_\mu_\nu}\right) + \delta\left(\sqrt{\abs g}\mathcal{L}_{\text{material}}(\chi, g)\right)\right]\\
    &= \int\dd^D x \left[ \frac{1}{2\kappa}\diffp{\sqrt{\abs g}g\indices{^\rho^\sigma}}{g\indices{^\mu^\nu}}\breve R\indices{_\rho_\sigma} + \diffp{\sqrt{\abs g}\mathcal{L}_{\text{mat}}}{g\indices{^\mu^\nu}}
    \right]\delta g\indices{^\mu^\nu} +\frac{1}{2\kappa}
    \sqrt{\abs g}g\indices{^\rho^\sigma}\delta\breve R\indices{_\rho_\sigma}\\
    &\hspace{1em}\Rightarrow \frac{1}{2\kappa}\diffp{\sqrt{\abs g}g\indices{^\rho^\sigma}}{g\indices{^\mu^\nu}}\breve R\indices{_\rho_\sigma} + \diffp{\sqrt{\abs g}\mathcal{L}_{\text{mat}}}{g\indices{^\mu^\nu}} = 0
\end{align*}
\begin{align*}
    \diffp{\sqrt{\abs g}g\indices{^\rho^\sigma}}{g\indices{^\mu^\nu}} &= \diffp{\sqrt{\abs g}}{g\indices{^\mu^\nu}}g\indices{^\rho^\sigma} + \sqrt{\abs g}\diffp{g\indices{^\rho^\sigma}}{g\indices{^\mu^\nu}} =  -\frac{1}{2}\sqrt{\abs g}g\indices{_\mu_\nu}g\indices{^\rho^\sigma} + \frac{1}{2}\sqrt{\abs g}(\delta^\rho_\mu\delta^\sigma_\nu + \delta^\rho_\nu\delta^\sigma_\mu)\\
    \diffp{\sqrt{\abs g}\mathcal{L}_{\text{mat}}}{g\indices{^\mu^\nu}} &=
    \sqrt{\abs g}\diffp{\mathcal{L}_{\text{mat}}}{g\indices{^\mu^\nu}} +  \diffp{\sqrt{\abs g}}{g\indices{^\mu^\nu}}\mathcal{L}_{\text{mat}}\\
    &= \sqrt{\abs g}\diffp{\mathcal{L}_{\text{mat}}}{g\indices{^\mu^\nu}} -\frac{1}{2}\sqrt{\abs g}{g\indices{_\mu_\nu}}\mathcal{L}_{\text{mat}} \equiv \frac{\sqrt{\abs g}}2 T\indices{_\mu_\nu}
\end{align*}
\begin{gather}
    \nonumber\left(\frac{1}{2}\sqrt{\abs g}(\delta^\rho_\mu\delta^\sigma_\nu + \delta^\rho_\nu\delta^\sigma_\mu) -\frac{1}{2}\sqrt{\abs g}g\indices{_\mu_\nu}g\indices{^\rho^\sigma}\right)\breve R\indices{_\rho_\sigma} =
    -\kappa\sqrt{\abs g} T\indices{_\mu_\nu}\\
    \breve G\indices{_\mu_\nu}\equiv\breve R\indices{_\rho_\sigma} -\frac{1}{2}g\indices{_\mu_\nu}\breve R = -8\pi T\indices{_\mu_\nu}
\end{gather}
onde $T\indices{_\mu_\nu}$ é o tensor enerxía-momento. A expresión (39) son as ecuacións de Einstein da relatividade xeral que notemos que están escritas en función da conexión de Levi-Civita, isto pódese xeneralizar a calquera conexión como comentaremos. No sucesivo suprimiremos a lagranxiana material por simplicidade.

\subsection{Formalismos de primeira orde} 
O formalismo de primeira orde pode ser enunciado de dúas formas, por unha banda está o formalismo que chamaremos \textit{de Cartan}, expresado mediante o vielbein; e pola outra está o formalismo de Palatini, que é a xeneralización do anterior pero considerando calquera conexión compatible coa métrica da forma que se sinala en (13). Esta segunda enunciación, que semella máis intuitiva vendo a forma anterior, leva aos mesmos resultados que se traballamos dende o fibrado tanxente; o formalismo de Palatini ten varias vantaxes como a inclusión natural da torsión ou unha ``doada'' xeneralización a $D>4$, o que quere dicir que a teorías de Lovelock son enunciables no formalismo de Palatini \cite{Lovelock_metricformulation}. Tamén se poderían demostrar os seguintes resultados\footnote{Non poñer ningún acento sobre os tensores supón que están construídos coa conexión de Palatini.}, remitimos a \cite{clave,Bernal_2017} para un estudo exaustivo,
\begin{align}
    G\indices{_\mu_\nu} &= R\indices{_{(\mu}_{\nu)}}-\frac{1}{2}R = -8\pi T\indices{_\mu_\nu}\hspace{2em}\text{Ecuacións de campo de Einstein-Palatini}\\
    R\indices{_{(\mu}_{\nu)}} &= \breve R\indices{_{\mu}_{\nu}}\ \Rightarrow\ R = \breve R \ \Rightarrow\ G\indices{_{\mu}_{\nu}} = \breve G\indices{_{\mu}_{\nu}}\\
    \Gamma\indices{^\mu_\alpha_\beta}& = \breve\Gamma\indices{^\mu_\alpha_\beta} -\frac{1}{D-1}\delta\indices{^\mu_\beta}T\indices{^\sigma_\sigma_\alpha}
    \hspace{2.206em}\text{Ecuación da conexión de Palatini}
\end{align}
e ademais, a nivel de traxectorias, as xeodésicas críticas e autoparalelas coinciden para calquera conexión de Palatini e incluso para dúas conexións distintas as traxectorias están relacionadas mediante un mero cambio de coordenadas, o que supón que as conexións de Palatini portan unha certa simetría física. Supúxose que a conexión de Palatini é compatible coa métrica ($g\indices{_\mu_\nu_{;\lambda}} = 0$), que para este caso é trivial recuperar o formalismo métrico mediante estes resultados.\\

Na sección anterior desenrolouse e probouse que do que se pode falar dende a linguaxe da variedade tamén pode ser expresado en linguaxe da tanxente. Deste xeito, é intuitivo pensar que a lagranxiana gravitatoria podése escribir en función do vielbein e da conexión de spin 
e ademais calquera formalismo lagranxiano escrito en 1-formas sempre é de primeira orde (non pode existir dependencia en derivadas superiores polo lema de Poincaré). Así, escribimos a acción de Einstein-Hilbert en 4D
\begin{align}
    \mathcal{S} = \int\varepsilon\indices{_a_b_c_d}(\alpha\mathcal{R}\indices{^a^b}e^ce^d + \beta e^ae^be^ce^d)
\end{align}
onde se suprimiron (e de aquí en adiante) os produtos wedge e o elemento $\varepsilon$ é o pseudotensor completamente antisimétrico de Levi-Civita, que é un pseudotensor invariante ($\dd\varepsilon=0$) \cite{CH_supergravities}. Se deducimos a ecuación de campo dende esta acción chegaríamos a unha ecuación de campo que só se correspondería ca anterior de Palatini se aplicamos uns postulados que se denominan postulado do vielbein e postulado de metricididade \cite{EinsteinsVielben}.

\section{Teorías de Lovelock}
A xeneralización natural a outras dimensións da acción de Einstein-Hilbert son as teorías de Lovelock, obxectivo final desta presentación, agora ben, estas non son tampouco as teorías máis xerais da relatividade. Como veremos, o teorema de Lovelock asume unha restrición bastante importante: o espazotempo é unha variedade libre de torsión. Unha teoría moito máis xeral sería considerar a torsión na nosa lagranxiana, isto sería un verdadeiro quebracabezas, tendo que traballar dende a formulación de Palatini \cite{Lovelock_metricformulation}, e requiriríamos de introducir as formas de Chern-Simmons \cite{CH_supergravities,Zanelli_2012} que escapa por moito as capacidades deste texto. 

\pagestyle{apaño1}
\subsection{Teorema de Lovelock}
% Enunciado do teorema
% Condiciones sobre as ecuacións do movemento
% Orden diferencial e conservación
Antes de enunciar o teorema, xa dixemos que se considera $\mathcal{T}^a=0$, polo que, se a conexión de lorentz non é independente do vielbein, teremos que $\omega\indices{^a^b}=\omega(e,\partial e)$. Logo, polo teorema de Poincaré, non poderá ser posible que aparezan derivadas superiores do vielbein contraído. A lagranxiana de Lovelock será do tipo $\mathcal{L}(\omega(e,\partial e),e)$ e, baixo a definición da curvatura 2-forma (33), podemos reescribila como $\mathcal{L}(\mathcal{R},e)$; por certa xeneralidade preferimos empregar $\mathcal{L}(\omega,e)$. Agora si, o teorema de Lovelock formulado entre 1970 e 1971 \cite{Lovelock-1970,Lovelock-1971} nunha enunciación compacta reza
\begin{adjustwidth}{1cm}{0.8cm}\vspace{0.4em}
    \textit{En ausencia de torsión, a acción máis xeral para a gravidade $\mathcal{S}(e,\omega)$ en D dimensións invariante baixo TLL é da forma}
\end{adjustwidth}
\begin{align}
    \mathcal{S}_D(e,\omega) = \kappa\int_\mathcal{M}\sum_{p=0}^{[D/2]}a_p\mathcal{L}_p^D
\end{align}
\begin{adjustwidth}{1cm}{0.8cm}
    \textit{onde $a_p$ é unha constante arbitraria e a lagranxiana vén descrita por}
\end{adjustwidth}
\begin{align}
    \mathcal{L}_p^D = \varepsilon_{a_1...a_D}\mathcal{R}\indices{^{a_1}^{a_2}}...\mathcal{R}\indices{^{a_{2p-1}}^{a_{2p}}}e\indices{^{a_{2p+1}}}...e\indices{^{a_D}}.
\end{align}
Coa definición dunha acción está claro o que temos que facer para desentrañar toda a súa información dinámica: aplicar o principio variacional. Se extremalizamos respecto de $e^a$ e $\omega\indices{^a^b}$ (44) temos unha expresión compacta da forma
\begin{align}
    \delta\mathcal{S}_D = \int\left(\delta e^a\mathcal{E}_a + \delta\omega\indices{^a^b}\mathcal{E}\indices{_a_b} \right) = 0
\end{align}
onde está claro que os termos $\mathcal{E}_a$ e 
$\mathcal{E}\indices{_a_b}$ débense anular simultaneamente para satisfacer que $\delta\mathcal{S}_D = 0$. Se fixeramos os cálculos varacionais con calma, semellante a como fixemos na sección 4.1., chegaríamos ás seguintes expresión para os termos $\mathcal{E}$
\begin{align}
    \mathcal{E}_a = \sum_{p=0}^{\left[\frac{D-1}{2}\right]}a^{(p)}(D-2p)\mathcal{E}_a^{(p)} = 0\ ;\\
    \mathcal{E}\indices{_a_b} = \sum_{p=1}^{\left[\frac{D-1}{2}\right]}a^{(p)}p(D-2p)\mathcal{E}_{ab}^{(p)} = 0
\end{align}
onde
\begin{align}
    \mathcal{E}_{a}^{(p)} &:=\ \varepsilon_{ac_1...c_{D-1}}\mathcal{R}\indices{^{c_1}^{c_2}}...\mathcal{R}\indices{^{c_{2p-1}}^{c_{2p}}}e\indices{^{c_{2p+1}}}...e\indices{^{c_{D-1}}}\ ;\\
    \mathcal{E}_{ab}^{(p)} &:=\ \varepsilon_{abc_1...c_{D-2}}\mathcal{R}\indices{^{c_1}^{c_2}}...\mathcal{R}\indices{^{c_{2p-3}}^{c_{2p-2}}}\mathcal{T}^{c_{2p-1}}e\indices{^{c_{2p}}}...e\indices{^{c_{D-2}}}.
\end{align}
A expresión (48) é trivialmente verificada xa que, por condicións do teorema de Lovelock, a torsión é nula, logo, o termo definido en (50) sempre se anula por ter un termo de torsión. Para (47) isto non pasa e teríamos que tratar cada caso concreto, no entanto, ambas condicións son polinomios de $\mathcal{R}_{ab}$ así que poderíamos atopar algunha combinación concreta baixo a que tanto (47) como (48) se anulen sen impoñer restricións sobre a torsión. Este resultado válenos para construír as ecuacións dinámicas $\forall \mathcal{S}_D$ de Lovelock.

% \subsection{Construción de invariantes de Lovelock}
% \textbf{DEBO CONSTRUÍR AS TEORÍAS EU}
% Densidades de Euler
% Polinomios de curvatura
% \subsection{Acción xeral de Lovelock}
%     expresión en D dimensións
%     coeficientes de acoplamento

\subsection{Acción de Lovelock en 2, 4 e 5 dimensións}
    % Termo cosmolóxico // métricas de de Sitter
    % Einstein–Hilbert
    % Lanczos-Gauss–Bonnet
    % Términos de orden superior
Rematamos o noso estudo afondando un pouco no cálculo concreto dalgunhas accións de Lovelock. Propoñerémonos calcular $\mathcal{S}_2$, $\mathcal{S}_3$, $\mathcal{S}_4$ e $\mathcal{S}_5$ en vista do teorema de Lovelock e as expresións (44) e (45). %D=2 (trivial) D=4 (EH + termo cosmolóxico) D=5 (gauss-bonnet coa respectiva ecuación de campo)
\paragraph{D=1+1.} O cálculo desta acción é inmediato, ten só dous termos, un relativo ao D-volume,$\int_\mathcal{M}\sqrt{\abs g}\dd^Dx$ do espazo e outro á curvatura
\begin{gather}
    \mathcal{L}_0^2 = \varepsilon_{ab}e^ae^b\hspace{1.5em}\text{e}\hspace{1.5em}\mathcal{L}_1^2 = \varepsilon_{ab}R^{ab}
\end{gather}\vspace{-2.5em}
\begin{align}
    \mathcal{S}_2 &= \kappa\int_\mathcal{M}\varepsilon_{ab}\left(
    a_1\mathcal{R}^{ab}+a_0e^ae^b
    \right)\\
    &= \kappa\int_\mathcal{M}\varepsilon_{ab}\left(
    \frac{a_1}2 \viel\biel R\indices{^\mu^\nu_{\rho\sigma}}\dd x^\rho\dd x^\sigma + a_0\viel\biel\dd x^\mu\dd x^\nu
    \right)\\
    &=\kappa\int_\mathcal{M}\sqrt{\abs g}\left(
    a_1R+2a_0
    \right)\dd^2x
\end{align}
Este caso non ten unha dinámica moi interesante, porén fixemos os cálculos un pouco máis exhaustivos; será o proceder que aplicaremos en xeral. Coidemos que basta con xogar cos índices introducindo a métrica para contraer o tensor de Riemann ao escalar de curvatura, ademais empregouse $\varepsilon_{a_1...a_D}e\indices{^{a_1}_{\mu_1}}...e\indices{^{a_D}_{\mu_D}}\dd x^{\mu_1}...\dd x^{\mu_D} = D!\sqrt{\abs g}\dd^Dx$. Coidemos que o termo cosmolóxico aparece $\forall D$ e correspóndese co termo $p=0$.
\vspace{-0.5em}
\paragraph{D=2+1 e D=3+1.} Aquí a acción redúcese á acción de Einstein-Hilbert cun termo D-volume engadido que coincidirá coa constante cosmolóxica. O caso $D=3$ é farto interesante (\cite{2+1grav0,2+1grav1}) e para $D=4$, que é a nosa gravidade usual, temos 
\begin{align}
    \mathcal{S}_4 &= \kappa\int_\mathcal{M}\varepsilon_{abcd}\left(
    a_2\mathcal{R}^{ab}\mathcal{R}^{cd} + a_1\mathcal{R}^{ab}e^ce^d + a_0e^ae^be^ce^d
    \right)\\ 
    &=\kappa\int_\mathcal{M}\sqrt{\abs g}\left[
    a_2\left( R^{\alpha\beta\gamma\delta}R_{\alpha\beta\gamma\delta} - 4R^{\alpha\beta}R_{\alpha\beta} + R^2 \right) + 2a_1R
    \right]\dd^4x + 24\kappa a_0\underbrace{\int_\mathcal{M}\sqrt{\abs g}\dd^4x}_{V_4}.
\end{align}
O termo que acompaña ao coeficiente $a_2$ é coñecido como densidade de Gauss-Bonnet-Lanczos (GBL), proporciona a primeira corrección non trivial a $\mathcal{S}_5$. En $D=4$ este termo non é relevante, é só topolóxico: non contribúe na dinámica \cite{Zanelli_2012}. Logo, a dinámica neste caso é a usual ecuación de campo de Einstein con termo cosmolóxico. 

\paragraph{D=5. Gravidade de Einstein-Gauss-Bonnet (EGB).}
%Un estudo pormenorizado do que sacamos os seguintes resultados é  \cite{Gauss_Bonnet_2022}, así metemos (1), dicimos que podemos recuperar EH con tal, e logo que tal termo é a cte cosmoloxica e escribimos (3), logo, facendo uso das ecuacións (47) e (48) podemos dar as ecuacións de campo (6) (7), aparece a primeira ecuación de campo con torsión non trivial, no entanto o teorema de lovelock fai que sexa cero. algun comentario chulo e xa
A gravidade EGB é un campo bastante activo actualmente e non imos afondar moito, para unha disertación máis pormenorizada \cite{Gauss_Bonnet_2022}. A primeira corrección que aparece é o termo GBL, así se desenrolamos a expresión (44) para $D=5$ teríamos a seguinte acción
\begin{align}
    \mathcal{S}_5 &= \kappa\int_\mathcal{M}\varepsilon_{abcdf}\left(
    a_2\mathcal{R}^{ab}\mathcal{R}^{cd}e^f + a_1\mathcal{R}^{ab}e^ce^d e^f + a_0e^ae^be^ce^de^f
    \right)\\
    &=\kappa\int_\mathcal{M}\sqrt{\abs g}\left[
    \alpha (GBL) + 4R - 8\Lambda_{5}
    \right]\dd^5x 
\end{align}
na que escollemos as constantes $a_1,a_0$ baixo unhas certas asuncións \cite{Gauss_Bonnet_2022}. Como para calquera outro caso aparece naturalmente o termo cosmolóxico, ademais notemos a equivalencia desta acción coa usual acción de Einstein-Hilbert (1) con tan só escoller $\alpha=0$. Logo, se queremos estudar a dinámica desta teoría da gravidade basta aplicar (47) e (48) coma en calquera outra acción e teríamos as ecuacións de campo
\begin{align}
    \varepsilon_{abcdf}\left(
    \alpha\mathcal{R}^{bc}\mathcal{R}^{df} + 2\mathcal{R}^{bc}e^d e^f - \frac{\Lambda_5}{3}e^be^ce^de^f
    \right) & = 0,\\
    \varepsilon_{abcdf}\left(
    2\alpha\mathcal{R}^{cd}\mathcal{T}^c + 2\mathcal{T}^c e^de^f
    \right) & = 0.
\end{align}
A ecuación (60) deixa explícita a torsión (algo que non atopabamos na formulación máis clásica), no entanto, o teorema de Lovelock esixe que a acción sexa libre torsión, porén (60) é trivial con $\mathcal{T}^a=0$. A ecuación (59) sería a correspondente ecuación de campo de Einstein para $5D$ con constante cosmolóxica e a corrección Gauss-Bonnet-Lanczos.
% Además, como lo demuestra
% el problema aún sin resolver de la constante cosmológica, no existe ningún argumento teórico
% que permita predecir sus valores mediante argumentos naturales de una manera que pueda compararse con
% las observaciones

En todo este desenrolo poderíamos ter incluído unha acción material como na sección 4.1. e teríamos que as ecuacións dinámicas (47) e (48) serían igualadas a $\sim T_{\mu\nu}$ e non a cero (o baleiro).
\pagestyle{apaño2}
\section{Conclusión}
No presente traballo presentouse unha detallada descrición matemática, tanto da xeometría riemanniana como do fibrado tanxente, que nos levaría a construír unha extensión natural da relatividade xeral en dimensións arbitrarias. 

Esta xeneralización é posible ao presentado teorema de Lovelock coa que construír accións do tipo (44), aínda que estas non son as teorías de gravidade máis xerais, pois son libres de torsión.
Os distintos termos de Lovelock constrúense como combinacións polinómicas da curvatura e inclúen como caso particular a acción de Einstein–Hilbert o que asegura certa consistencia.
Finalmente, observouse que en catro dimensións a teoría redúcese esencialmente á relatividade xeral, mentres que para dimensións superiores aparecen novas contribucións dinámicas como a densidade GBL.
